% GDZ PPN236018647,native572lower–573upper;HKUST573–574.
\NumberedItem{99}
\indent In principiis Geometriae egregios progressus fecimus.
\par\noindent\hfill Br[unsvigae, 1799] Sept.

\medskip\begingroup\small
\indent Für diese Arbeiten über die ersten Gründe der Geometrie siehe den Brief von \textsc{Gauss} an \textsc{Wolfgang Bolyai} vom 16. Dezember 1799, Werke VIII, S. 159. Wahrscheinlich stammt auch der Beweis des Lehr-
\par\endgroup

\clearpage\SourceLeaf{573}\ArithmeticHead{544}{L}{TAGEBUCH MIT ERLÄUTERUNGEN.}
\begingroup\small
\noindent satzes vom Flächeninhalt des Dreiecks, den \textsc{Gauss} in dem Briefe an \textsc{Wolfgang Bolyai} vom 6. März 1832 mitteilt (Werke VIII, S. 221), aus dieser frühen Zeit; vergl. den Aufsatz »\textsc{Gauss} als Geometer«, Werke X\,2.
\par\noindent\hfill\textsc{Stäckel.}
\par\endgroup
